\heading{高等数学A（II）-23春-期中}
\noteinfo{题目由用户提供，解答按本仓库现有版式整理。}

\begin{question}[指出积分类型并计算]

\begin{enumerate}[label=(\arabic*),leftmargin=2.6em,itemsep=0.25em,topsep=0.2em]
\item $\int_{[0,1]}x\,dx$；
\item $\oint_{\partial[0,1]^2}xy\,ds$；
\item $\oiint_{\partial[0,1]^3}xyz\,dS$；
\item $\iint_{[0,1]^2}xy\,dx\,dy$；
\item $\oint_{\partial[0,1]^2}2xy\,dx+(x^2+y^2)\,dy$；
\item $\iiint_{[0,1]^3}x^6y^{16}z^{16}\,dV$；
\item
\[
\oiint_{\partial[0,1]^3}\left(\frac{x}{2}+z^3\sin y^2\right)\,dy\,dz
+\left(\frac{y}{3}+e^{x\cos z}\right)\,dz\,dx
+\left(\frac{z}{6}+\arctan(xy)\right)\,dx\,dy;
\]
\item
\[
\oint_{\Gamma}x\,dx+y\,dy+z\,dz,
\]
其中 $\Gamma$ 为折线段
\[
(0,0,0)\to(1,0,0)\to(1,1,0)\to(1,1,1)\to(0,1,1)\to(0,1,0)\to(0,0,0).
\]
\end{enumerate}

\begin{solution}
\begin{enumerate}[label=(\arabic*),leftmargin=2.6em,itemsep=0.25em,topsep=0.2em]
\item 这是\textbf{一元定积分}，
\[
\int_0^1 x\,dx=\left.\frac{x^2}{2}\right|_0^1=\frac12.
\]
故
\ansmath{\int_{[0,1]}x\,dx=\frac12.}

\item 这是\textbf{第一类曲线积分}。单位正方形四条边中，位于 $x=0$ 或 $y=0$ 的两条边贡献为 $0$；位于 $x=1$ 与 $y=1$ 的两条边分别给出
\[
\int_0^1 y\,dy=\frac12,
\qquad
\int_0^1 x\,dx=\frac12.
\]
故
\ansmath{\oint_{\partial[0,1]^2}xy\,ds=1.}

\item 这是\textbf{第一类曲面积分}。单位立方体六个面中，$x=0$、$y=0$、$z=0$ 三个面上的被积函数为 $0$；在 $x=1$、$y=1$、$z=1$ 三个面上分别有
\[
\iint yz\,dy\,dz=\frac14,\qquad
\iint xz\,dx\,dz=\frac14,\qquad
\iint xy\,dx\,dy=\frac14.
\]
因此
\ansmath{\oiint_{\partial[0,1]^3}xyz\,dS=\frac34.}

\item 这是\textbf{二重积分}，
\[
\iint_{[0,1]^2}xy\,dx\,dy
=\left(\int_0^1 x\,dx\right)\left(\int_0^1 y\,dy\right)
=\frac14.
\]
故
\ansmath{\iint_{[0,1]^2}xy\,dx\,dy=\frac14.}

\item 这是\textbf{第二类曲线积分}。记
\[
P(x,y)=2xy,\qquad Q(x,y)=x^2+y^2.
\]
由 Green 公式，
\[
\oint_{\partial[0,1]^2}P\,dx+Q\,dy
=\iint_{[0,1]^2}\left(\frac{\partial Q}{\partial x}-\frac{\partial P}{\partial y}\right)\,dx\,dy
=\iint_{[0,1]^2}(2x-2x)\,dx\,dy=0.
\]
故
\ansmath{\oint_{\partial[0,1]^2}2xy\,dx+(x^2+y^2)\,dy=0.}

\item 这是\textbf{三重积分}，
\[
\iiint_{[0,1]^3}x^6y^{16}z^{16}\,dV
=\left(\int_0^1 x^6\,dx\right)\left(\int_0^1 y^{16}\,dy\right)\left(\int_0^1 z^{16}\,dz\right)
=\frac{1}{7\cdot17\cdot17}.
\]
注意
\[
7\cdot17\cdot17=2023,
\]
故
\ansmath{\iiint_{[0,1]^3}x^6y^{16}z^{16}\,dV=\frac{1}{2023}.}

\item 这是\textbf{第二类曲面积分}。记
\[
P=\frac{x}{2}+z^3\sin y^2,\qquad
Q=\frac{y}{3}+e^{x\cos z},\qquad
R=\frac{z}{6}+\arctan(xy).
\]
由 Gauss 公式，
\[
\oiint_{\partial[0,1]^3}P\,dy\,dz+Q\,dz\,dx+R\,dx\,dy
=\iiint_{[0,1]^3}\left(\frac{\partial P}{\partial x}+\frac{\partial Q}{\partial y}+\frac{\partial R}{\partial z}\right)\,dV.
\]
而
\[
\frac{\partial P}{\partial x}=\frac12,\qquad
\frac{\partial Q}{\partial y}=\frac13,\qquad
\frac{\partial R}{\partial z}=\frac16,
\]
故散度恒为 $1$。因此
\ansmath{\oiint_{\partial[0,1]^3}\cdots=1.}

\item 这是\textbf{第二类曲线积分}。注意
\[
x\,dx+y\,dy+z\,dz=d\!\left(\frac{x^2+y^2+z^2}{2}\right).
\]
故它是全微分沿闭曲线的积分，从而
\ansmath{\oint_{\Gamma}x\,dx+y\,dy+z\,dz=0.}
\end{enumerate}
\end{solution}
\end{question}

\begin{question}[二重积分]

计算
\[
I=\iint_{-1\le x\le1,\ 0\le y\le1}|y-x|\,dx\,dy.
\]

\begin{solution}
按直线 $x=y$ 分区域，有
\[
I=\int_0^1\left(\int_{-1}^y (y-x)\,dx+\int_y^1 (x-y)\,dx\right)\,dy.
\]
先算内层积分：
\[
\int_{-1}^y (y-x)\,dx=\frac{y^2}{2}+y+\frac12,
\]
\[
\int_y^1 (x-y)\,dx=\frac12-y+\frac{y^2}{2}.
\]
故
\[
I=\int_0^1 (y^2+1)\,dy
=\left.\left(\frac{y^3}{3}+y\right)\right|_0^1
=\frac43.
\]
因此
\ansmath{I=\frac43.}
\end{solution}
\end{question}

\begin{question}[求体积]

求闭曲面
\[
\left(x^2+\frac{y^2}{4}+\frac{z^2}{9}\right)^2\le x
\]
围成区域的体积。

\begin{solution}
作线性代换
\[
u=x,\qquad v=\frac y2,\qquad w=\frac z3.
\]
则
\[
x=u,\qquad y=2v,\qquad z=3w,
\]
其 Jacobian 为
\[
\left|\frac{\partial(x,y,z)}{\partial(u,v,w)}\right|=6.
\]
原区域化为
\[
(u^2+v^2+w^2)^2\le u.
\]
在以 $u$ 轴为极轴的球坐标中，
\[
u=r\cos\varphi,\qquad
v=r\sin\varphi\cos\theta,\qquad
w=r\sin\varphi\sin\theta.
\]
于是区域满足
\[
r^4\le r\cos\varphi,
\]
故
\[
0\le \theta\le 2\pi,\qquad
0\le \varphi\le \frac{\pi}{2},\qquad
0\le r\le (\cos\varphi)^{1/3}.
\]
因此在 $(u,v,w)$ 空间中的体积为
\[
V_{uvw}
=\int_0^{2\pi}\!\int_0^{\pi/2}\!\int_0^{(\cos\varphi)^{1/3}}r^2\sin\varphi\,dr\,d\varphi\,d\theta
=2\pi\cdot\frac13\int_0^{\pi/2}\cos\varphi\sin\varphi\,d\varphi
=\frac{\pi}{3}.
\]
所以原区域体积
\[
V=6V_{uvw}=2\pi.
\]
故
\ansmath{V=2\pi.}
\end{solution}
\end{question}

\begin{question}[曲面积分]

计算
\[
I=\oiint_{x^2+y^2+z^2=1}\frac{x\,dy\,dz+y\,dz\,dx+z\,dx\,dy}{(x^2+4y^2+9z^2)^{3/2}}.
\]

\begin{solution}
作代换
\[
u=x,\qquad v=2y,\qquad w=3z.
\]
则
\[
dy=\frac12\,dv,\qquad dz=\frac13\,dw,\qquad dx=du,
\]
并且
\[
x\,dy\,dz+y\,dz\,dx+z\,dx\,dy
=\frac16\bigl(u\,dv\,dw+v\,dw\,du+w\,du\,dv\bigr).
\]
故
\[
I=\frac16\oiint_{\Sigma}\frac{u\,dv\,dw+v\,dw\,du+w\,du\,dv}{(u^2+v^2+w^2)^{3/2}},
\]
其中 $\Sigma$ 为椭球面
\[
u^2+\frac{v^2}{4}+\frac{w^2}{9}=1.
\]

记向量场
\[
\mathbf F(u,v,w)=\frac{(u,v,w)}{(u^2+v^2+w^2)^{3/2}}.
\]
它在原点外满足 $\nabla\cdot\mathbf F=0$，故其穿过任意包围原点的闭曲面的通量都相同。于是可将 $\Sigma$ 连续变形为单位球面 $u^2+v^2+w^2=1$。在单位球面上，
\[
\mathbf F=(u,v,w),\qquad \mathbf n=(u,v,w),
\]
故
\[
\mathbf F\cdot\mathbf n=1.
\]
所以通量等于球面积 $4\pi$，从而
\[
I=\frac16\cdot4\pi=\frac{2\pi}{3}.
\]
故
\ansmath{I=\frac{2\pi}{3}.}
\end{solution}
\end{question}

\begin{question}[三重积分换序]

设 $f$ 在 $[0,1]$ 上可积，求
\[
I=\int_0^1\int_0^x\int_0^y f(z)\,dz\,dy\,dx.
\]

\begin{solution}
积分区域为
\[
0\le z\le y\le x\le1.
\]
交换积分次序，先对 $x,y$ 积分，得
\[
I=\int_0^1 f(z)\left(\int_z^1\int_y^1 dx\,dy\right)\,dz.
\]
而
\[
\int_y^1 dx=1-y,
\]
故
\[
\int_z^1(1-y)\,dy=\frac{(1-z)^2}{2}.
\]
于是
\[
I=\frac12\int_0^1(1-z)^2f(z)\,dz.
\]
因此
\ansmath{I=\frac12\int_0^1(1-z)^2f(z)\,dz.}
\end{solution}
\end{question}

\begin{question}[极限]

设
\[
F(t)=\iiint_{x^2+y^2+z^2\le t^2}f\!\left(x^2+y^2+z^2\right)\,dV,
\]
其中 $f$ 连续，在 $0$ 处可导，且 $f(0)=0,\ f'(0)=10$。求
\[
\lim_{t\to0^+}\frac{F(t)}{t^5}.
\]

\begin{solution}
用球坐标，令 $r=\sqrt{x^2+y^2+z^2}$，则积分区域为 $0\le r\le t$，体积元 $dV=r^2\sin\varphi drd\varphi d\theta$，于是
\[
F(t)=4\pi\int_0^t r^2f\!\left(r^2\right)\,dr.
\]
由 $f$ 在 $0$ 处可导且 $f(0)=0$，可得
\[
f(u)=f(0)+f'(0)u+o(u)=10u+o(u)\qquad (u\to0^+).
\]
代入 $u=r^2$，得
\[
f(r^2)=10r^2+o(r^2)\quad(r\to0^+).
\]
因此
\[
F(t)=4\pi\int_0^t r^2\bigl(10r^2+o(r^2)\bigr)\,dr
=4\pi\int_0^t\bigl(10r^4+o(r^4)\bigr)\,dr
=8\pi t^5+o(t^5).
\]
于是
\[
\frac{F(t)}{t^5}=\frac{8\pi t^5+o(t^5)}{t^5}=8\pi+o(1)\xrightarrow[t\to0^+]{}8\pi.
\]
故按题面所求，
\ansmath{\lim_{t\to0^+}\frac{F(t)}{t^5}=8\pi.}
\end{solution}
\end{question}

\begin{question}[圆域与方域的广义积分]

设
\[
I_r=\iint_{x^2+y^2\le r^2}f(x,y)\,dx\,dy,
\qquad
J_\rho=\iint_{|x|,|y|\le\rho}f(x,y)\,dx\,dy.
\]
证明：若 $\lim\limits_{r\to\infty}I_r$ 与 $\lim\limits_{\rho\to\infty}J_\rho$ 之一存在且有限，则另一者也存在且有限，且相等。

\begin{solution}
记
\[
D_r=\{(x,y):x^2+y^2\le r^2\},
\qquad
Q_\rho=\{(x,y):|x|\le\rho,\ |y|\le\rho\}.
\]
显然有集合包含关系
\[
D_\rho\subset Q_\rho\subset D_{\sqrt2\,\rho}.
\]

先设
\[
\lim_{r\to\infty}I_r=L
\]
存在且有限。由于 $I_r$ 收敛，所以对任意 $\varepsilon>0$，存在 $R>0$，使得当 $r_2>r_1>R$ 时，
\[
|I_{r_2}-I_{r_1}|<\varepsilon.
\]
取 $\rho>R$，由
\[
D_\rho\subset Q_\rho\subset D_{\sqrt2\,\rho}
\]
可得
\[
J_\rho-I_\rho=\iint_{Q_\rho\setminus D_\rho}f(x,y)\,dx\,dy,
\]
而集合 $Q_\rho\setminus D_\rho$ 包含在环域 $D_{\sqrt2\,\rho}\setminus D_\rho$ 中，因此
\[
|J_\rho-I_\rho|
\le \left|I_{\sqrt2\,\rho}-I_\rho\right|
<\varepsilon
\]
（按广义积分的 Cauchy 性理解即可）。于是
\[
J_\rho-I_\rho\to0.
\]
又 $I_\rho\to L$，故
\[
J_\rho\to L.
\]

反过来，若
\[
\lim_{\rho\to\infty}J_\rho=L
\]
存在且有限，由
\[
Q_{r/\sqrt2}\subset D_r\subset Q_r
\]
同理可得
\[
I_r-J_{r/\sqrt2}\to0.
\]
再由 $J_{r/\sqrt2}\to L$，便知
\[
I_r\to L.
\]

综上，只要其中一个极限存在且有限，另一个也存在且有限，并且
\ansmath{\lim_{r\to\infty}I_r=\lim_{\rho\to\infty}J_\rho.}
\end{solution}
\end{question}

